\documentclass[
    language=english,
    doctype=dissertation,
    institution=none,
    titlestyle=book,
]{../../omnilatex}

\RequirePackage{config/document-settings}
\addbibresource{../../bib/bibliography.bib}

\begin{document}

\maketitle

\frontmatter
\chapter{Abstract}

This dissertation investigates the design and optimization of adaptive streaming algorithms for real-time analysis of large-scale data streams. We propose a novel framework, called \emph{FluxStream}, that dynamically adjusts its resource allocation based on workload characteristics observed over sliding windows. Our contributions include: (1)~a formal model for adaptive stream processing that generalizes prior work on fixed-budget algorithms, (2)~an online learning component that converges to near-optimal resource allocation within logarithmic time, and (3)~an empirical evaluation on both synthetic benchmarks and real-world traces from network monitoring and financial data feeds.

Experimental results demonstrate that FluxStream reduces processing latency by 37\% on average compared to state-of-the-art static schedulers while maintaining 99.97\% answer accuracy for frequency estimation and quantile queries. The system scales linearly to throughput rates exceeding $10^6$ events per second on commodity hardware.

\tableofcontents
\listoffigures
\listoftables

\mainmatter

\chapter{Introduction}\label{ch:intro}

\section{Motivation}

The proliferation of sensors, social media, and online transaction systems has generated an unprecedented volume of continuous data. Traditional batch-processing paradigms are fundamentally mismatched to the low-latency requirements of applications such as intrusion detection, algorithmic trading, and real-time recommendation engines. Stream processing systems must analyze data incrementally, maintain summaries under strict memory constraints, and deliver answers with provable accuracy guarantees.

Despite significant advances in streaming algorithms over the past two decades, most existing systems employ static resource allocation strategies that are configured at deployment time and remain fixed throughout execution. This rigid approach leads to either over-provisioning (wasting resources during low-traffic periods) or under-provisioning (causing latency spikes and accuracy degradation during surges).

\section{Problem Statement}

We address the following research questions:

\begin{enumerate}
    \item How can a streaming system \emph{autonomously} adjust its computational budget across multiple concurrent queries to minimize end-to-end latency subject to accuracy constraints?
    \item What theoretical guarantees can be provided on the convergence rate and steady-state optimality of such an adaptive system?
    \item How does the proposed approach compare empirically with existing static and manually tuned systems across diverse workloads?
\end{enumerate}

\section{Contributions}

The primary contributions of this dissertation are:

\begin{enumerate}
    \item A \emph{unified cost model} that captures the trade-off between approximation error, memory usage, and processing latency for common stream queries (heavy hitters, quantiles, distinct counting).
    \item The \emph{FluxStream} adaptive framework, which combines bandit-based online learning with a lightweight performance monitor to dynamically reassign memory and CPU shares among queries.
    \item A comprehensive experimental evaluation spanning four synthetic workloads and two real-world datasets, demonstrating consistent improvements in latency and robustness.
\end{enumerate}

\section{Dissertation Outline}

The remainder of this dissertation is organized as follows. Chapter~\ref{ch:related} surveys related work in streaming algorithms and adaptive resource management. Chapter~\ref{ch:methodology} presents the FluxStream framework, including the cost model, the online learning algorithm, and the system architecture. Chapter~\ref{ch:results} reports experimental results. Chapter~\ref{ch:discussion} discusses implications, limitations, and directions for future work. Chapter~\ref{ch:conclusion} concludes with a summary of findings.

\chapter{Related Work}\label{ch:related}

\section{Streaming Algorithms}

The field of data stream algorithms was formalized in the seminal work of Alon, Matias, and Szegedy (1996), who introduced the frequency-moments problem and demonstrated that polynomial-space computation could be replaced by logarithmic-space sketches. Subsequent work established algorithms for heavy-hitter detection, quantile estimation, set membership, and distinct element counting. Cormode and Garofalakis (2005) provided an influential survey of sketch-based approaches, while Aggarwal (2007) edited a comprehensive treatment of the theoretical foundations.

These algorithms typically assume a \emph{fixed} resource budget determined at compile time. Adaptive variants exist---for example, the ``timestampl'' framework of Braverman and Ostrovsky (2010) allows algorithms to choose precision dynamically---but they require prior knowledge of the data distribution.

\section{Adaptive Resource Management}

In the broader systems literature, adaptive resource management has been studied in the context of operating systems (process scheduling with feedback control), databases (query plan adaptation via mid-query reoptimization), and cloud computing (elastic scaling of virtual machines). The most relevant line of work is the stream-processing literature on load shedding and approximate query processing.

Load-shedding techniques, introduced by Babcock et al.\ (2004), selectively drop tuples during overload to keep latency bounded. However, load shedding fundamentally sacrifices data completeness. Our approach instead redistributes computational resources to maintain approximate answers over the \emph{entire} stream.

\section{Online Learning in Systems}

Multi-armed bandit algorithms have been applied to system optimization in settings such as web server configuration, energy management in data centers, and compiler flag tuning. Our use of contextual bandits differs from prior work in that the context signal (observed stream statistics) changes continuously, requiring a non-stationary bandit formulation.

\chapter{Methodology}\label{ch:methodology}

\section{System Model and Cost Function}

We model a stream processing system as a set $Q = \{q_1, q_2, \dots, q_k\}$ of concurrent queries sharing a bounded memory pool of size $M$ and a CPU with total capacity $C$. Each query $q_i$ is associated with a sketch data structure $S_i(m_i)$ parameterized by its allocated memory $m_i$, where $\sum_{i=1}^{k} m_i \leq M$. The error of query $q_i$ is a monotonically decreasing function $\epsilon_i(m_i)$ of its memory allocation, and the per-element processing time is $\tau_i(m_i)$.

The system objective is to minimize the weighted average latency
\[
L = \sum_{i=1}^{k} w_i \cdot \tau_i(m_i)
\]
subject to $\epsilon_i(m_i) \leq \bar{\epsilon}_i$ for all $i$ and $\sum_{i} m_i \leq M$, where $w_i$ are user-specified priority weights and $\bar{\epsilon}_i$ are accuracy targets.

\section{Online Learning Algorithm}

We formulate the resource allocation problem as a contextual bandit. At each time step $t$, the system observes a context vector $\mathbf{x}_t \in \mathbb{R}^d$ encoding recent stream statistics (arrival rate, skewness, domain cardinality) and selects an allocation vector $\mathbf{m}_t$. The reward is the negative latency:
\[
r_t = -L(\mathbf{m}_t, \mathbf{x}_t).
\]

We employ a variant of Thompson sampling with a Gaussian reward model. Each arm corresponds to a discretized allocation profile, and the posterior over expected rewards is updated incrementally as feedback arrives from the performance monitor. To handle non-stationarity, we apply a sliding-window restart: the posterior is reset every $W$ time steps, where $W$ is chosen to be $O(\log k)$ times the mixing time of the stream.

\section{System Architecture}

FluxStream is implemented as a middleware layer between the input data source and the query execution engine. It consists of three components:

\begin{enumerate}
    \item \textbf{Monitor}: Collects per-query latency and error statistics at millisecond granularity using lightweight counters.
    \item \textbf{Learner}: Runs the bandit algorithm in a dedicated thread, recomputing the allocation every $T_{\text{adapt}}$ seconds ($T_{\text{adapt}} = 100\,\text{ms}$ in our experiments).
    \item \textbf{Executor}: Applies the new allocation by migrating sketch data structures between memory pools and adjusting CPU time slices via the OS scheduler.
\end{enumerate}

\section{Experimental Setup}

We evaluate FluxStream on two categories of workloads:

\begin{itemize}
    \item \textbf{Synthetic benchmarks}: Zipf-distributed integer streams with varying skew parameters ($z \in \{0.5, 1.0, 1.5, 2.0\}$) and throughput rates from $10^4$ to $10^6$ events per second. Queries include Count-Min sketch frequency estimation, Greenwald--Khanna quantile computation, and HyperLogLog distinct counting.
    \item \textbf{Real-world traces}: (1)~NetFlow records from a campus network gateway (7 days, 2.3 billion records), and (2)~NASDAQ ITCH feed data from one trading day (1.1 billion messages).
\end{itemize}

All experiments were conducted on a machine with dual Intel Xeon E5-2680 v4 CPUs (56 logical cores), 256\,GB RAM, running Ubuntu 22.04 LTS. FluxStream is implemented in C++17 and compiled with GCC 12.2 at \texttt{-O2}.

\chapter{Results}\label{ch:results}

\section{Latency Reduction}

Figure~\ref{fig:latency} (hypothetical) shows the 99th-percentile tail latency for FluxStream and three baselines: uniform allocation, proportionally fair allocation, and a manually tuned static configuration. Across all four synthetic workloads, FluxStream achieves the lowest latency, with improvements ranging from 28\% (Zipf $z = 0.5$) to 49\% (Zipf $z = 2.0$).

On the NetFlow trace, FluxStream reduces tail latency by 37\% compared to the manually tuned baseline. On the NASDAQ trace, the improvement is 41\%. The larger gains on skewed data are consistent with the hypothesis that adaptive allocation is most beneficial when query difficulty varies over time.

\section{Accuracy Maintenance}

Throughout all experiments, FluxStream maintains per-query accuracy within 0.5 percentage points of the user-specified targets ($\bar{\epsilon}_i = 1\%$ for frequency estimation, $0.5\%$ for quantile queries). No accuracy violation exceeds $1.2\%$, and such violations occur only during the first $200\,\text{ms}$ after a sharp workload transition.

\section{Convergence Speed}

The bandit learner converges to within 5\% of the optimal allocation in a median of $O(\log k)$ adaptation steps. For $k = 8$ queries, convergence typically occurs within 1--2 seconds. This rapid convergence ensures that the system responds effectively to bursty traffic patterns.

\section{Overhead Analysis}

The Monitor, Learner, and Executor components collectively consume less than 3\% of total CPU time and less than 1\% of the memory budget. The primary overhead comes from sketch migration during reallocation, which costs $O(m_i)$ memory copies per query. In practice, this overhead is amortized because reallocations occur infrequently (once per $T_{\text{adapt}}$ interval).

\chapter{Discussion}\label{ch:discussion}

\section{Implications}

The results demonstrate that online learning is a viable and practical approach to adaptive resource management in stream processing systems. By framing the allocation problem as a contextual bandit, FluxStream achieves near-optimal performance without requiring any prior knowledge of the workload distribution. This is particularly valuable in production environments where workloads are heterogeneous, non-stationary, and difficult to characterize a priori.

\section{Limitations}

Several limitations warrant discussion. First, our cost model assumes independent queries, whereas in practice queries may share intermediate results or operate on overlapping data subsets. Extending FluxStream to handle query interdependencies is an important direction for future work. Second, the discretized arm space limits the granularity of allocation adjustments. A continuous-action variant using policy gradient methods could potentially yield finer-grained control.

\section{Future Work}

We identify three promising avenues for future research:

\begin{enumerate}
    \item \textbf{Multi-objective optimization}: Extending the framework to simultaneously optimize latency, accuracy, and energy consumption using Pareto-optimal bandit policies.
    \item \textbf{Distributed settings}: Scaling FluxStream across multiple machines in a cluster, where each node runs a local bandit agent and periodically exchanges allocation hints with peers.
    \item \textbf{Learned cost models}: Replacing the hand-crafted error and latency functions $\epsilon_i(m_i)$ and $\tau_i(m_i)$ with neural-network surrogates trained on profiling data, enabling the system to adapt to novel sketch implementations without manual modeling.
\end{enumerate}

\chapter{Conclusion}\label{ch:conclusion}

This dissertation presented FluxStream, an adaptive framework for resource management in stream processing systems. The key insight is that workload characteristics vary sufficiently over time to justify online, data-driven resource allocation. By combining a principled cost model with contextual bandit learning, FluxStream achieves substantial latency reductions while maintaining strict accuracy guarantees.

The empirical evaluation on both synthetic and real-world data demonstrates that FluxStream outperforms static allocation strategies across a wide range of conditions. The system is lightweight, converges quickly, and requires no manual tuning beyond the specification of accuracy targets.

We believe that the adaptive paradigm explored here has broad applicability beyond stream processing. Any system that manages shared resources under dynamic workloads---from database query optimizers to operating system schedulers---could benefit from similar online-learning techniques.

\appendix
\chapter{Supplementary Material}\label{app:supplementary}

\section{Proof of Theorem 1}

\begin{theorem}
Under the non-stationary contextual bandit formulation with sliding-window restart, the expected cumulative regret of FluxStream over $T$ time steps is bounded by $O(\sqrt{dT \log k})$, where $d$ is the context dimension and $k$ is the number of queries.
\end{theorem}

\begin{proof}
We decompose the regret into per-window segments. Each window has length $W = O(\log k \cdot \tau_{\text{mix}})$, where $\tau_{\text{mix}}$ is the mixing time of the stream. Within a single window, the environment is approximately stationary, so standard Thompson sampling analysis yields per-window regret of $O(\sqrt{dW \log k})$.

The number of windows is $T/W$. Summing over all windows and applying Jensen's inequality:
\[
R_T \leq \frac{T}{W} \cdot O(\sqrt{dW \log k}) = O\!\left(\sqrt{\frac{dT^2 \log k}{W}}\right) = O(\sqrt{dT \log k}).
\]
\end{proof}

\section{Additional Experimental Data}

Table~\ref{tab:detailed} presents detailed per-query latency and accuracy numbers for the NetFlow trace experiment.

\begin{table}[htbp]
\centering
\caption{Per-query performance on NetFlow trace}\label{tab:detailed}
\begin{tabular}{lccccc}
\toprule
\textbf{Query} & \textbf{Type} & \textbf{Avg Lat (ms)} & \textbf{P99 Lat (ms)} & \textbf{Avg Err (\%)} & \textbf{Memory (MB)} \\
\midrule
Q1 & Frequency & 0.12 & 0.34 & 0.41 & 64 \\
Q2 & Frequency & 0.09 & 0.28 & 0.52 & 32 \\
Q3 & Quantile  & 0.18 & 0.51 & 0.63 & 128 \\
Q4 & Quantile  & 0.15 & 0.42 & 0.58 & 64 \\
Q5 & Distinct  & 0.22 & 0.67 & 0.71 & 256 \\
Q6 & Distinct  & 0.19 & 0.58 & 0.44 & 128 \\
\bottomrule
\end{tabular}
\end{table}

\printbibliography

\end{document}
